\section{Preserve the adjustment question across the DAGitty wrapper}
\label{exec18:gp-graph}

A bank analyst who supplies a causal diagram still needs to specify which
intervention and outcome the diagram is meant to relate. A graph editor, a
test of an adjustment set, an enumeration of such sets and an estimator of an
effect answer different questions. Our proposed use of DAGitty is deliberately
narrow: a supplied observed directed acyclic graph, one exposure, one outcome
and a total-effect adjustment query. The proposed first implementation would
call a pinned original-language R wrapper and preserve its graph and query
semantics. This is a proposal selection, not an installed adapter or a claim
that the bank's supplied arrows are true.
\citep{exec18_gp_dagitty}

The selected repository commit is dated 6 December 2024; its R description
declares version 0.3-4 and date 7 December 2023. Those dates identify the
inspected historical source and do not establish present support. The README
describes the shared JavaScript core and the R/V8 boundary. The R adjustment
entry points check that exposure and outcome variables are supplied and stop
with an error if either is absent. The inspected JavaScript Boolean check
instead returns false in that case. An application that interprets every false
as a scientific finding that adjustment is impossible would turn a missing
question into a causal conclusion. The adapter must validate the question
before interpreting the upstream value. \citep{exec18_gp_dagitty}

For a complete supported query, three distinct outputs remain necessary.
A valid empty covariate set means no adjustment is required; an empty
collection of valid sets means that this adjustment search supplied none;
an incomplete search or invalid request remains unresolved. In particular,
\(\{\varnothing\}\) and \(\varnothing\) have different meanings. The first
contains one set and the second contains none. Even an exhaustive failure
to find an adjustment set does not exclude identification through a method
outside covariate adjustment. The documented \texttt{max.results} argument
caps minimal-set enumeration. A capped listing cannot be reported as a list
of all valid choices, and an interrupted or unqualified computation cannot
be promoted to a proof of nonidentification. These distinctions belong in
the returned result and its completeness field. \citep{exec18_gp_dagitty}

\subsection{Three banking laws expose different conversion errors}
\label{exec18:gp-graph-laws}

Consider three separate project-authored structural laws. Their probabilities
are stipulated; they are mathematical references for a future adapter, not
observed customers or executed DAGitty examples. In the first, \(A\) is a
binary intervention, \(Y\) a binary favorable repayment outcome and \(Z\) a
baseline readiness attribute. The graph is \(Z\to A\), \(Z\to Y\),
\(A\to Y\), with mutually independent structural randomizers and
\begin{equation}
\label{exec18:gp-confounding}
 \begin{aligned}
 \Pr(Z=1)&=\frac12, &
 \Pr(A=1\mid Z=z)&=\frac{1+2z}{4},\\
 \Pr(Y=1\mid A=a,Z=z)&=\frac{1+2a+z}{8}.
 \end{aligned}
\end{equation}
Every relevant treatment probability is positive. Under this declared graph
and law, standardizing over the population distribution of \(Z\) gives
\(\Pr\{Y(0)=1\}=3/16\) and \(\Pr\{Y(1)=1\}=7/16\). Their difference is
\(1/4\). The unadjusted observed risks are instead \(5/32\) and \(15/32\),
whose difference is \(5/16\). Here the supplied baseline variable is needed
for the adjustment interpretation. A wrapper that drops the selected set
would change the target calculation. Exact enumeration of all eight joint
cells reproduces both pairs. \citep{exec18_gp_graph_reference}

The second law shows why adding variables is not a universal remedy. Let
\(A\) and an unobserved resilience variable \(U\) be independent
Bernoulli\((1/2)\) variables, let \(Y=U\), and let documentation completion
\(C\) satisfy
\begin{equation}
\label{exec18:gp-collider}
 \Pr(C=1\mid A=a,U=u)=\frac{1+2a+2u}{8},
 \qquad A\to C\leftarrow U\to Y.
\end{equation}
There is no directed path from \(A\) to \(Y\), so the total effect is zero.
The empty adjustment set is valid. Nevertheless, the observed repayment
contrasts within \(C=0\) and \(C=1\) are respectively \(-1/24\) and
\(-1/8\). Standardizing these contrasts over the marginal distribution
of \(C\) gives \(-7/96\). All eight \((A,U,C)\) cells have positive probability;
the discrepancy is caused by conditioning on the collider, rather than by
a missing-support cell. A business workflow that automatically adds a
post-intervention documentation field can therefore introduce a spurious
adverse relationship. The graph and the information clock must survive
conversion together. \citep{exec18_gp_graph_reference}

In the third law, \(A\) is randomized, \(M\) is subsequent engagement,
and the graph contains \(A\to M\to Y\) and \(A\to Y\). Set
\(\Pr(M=1\mid A=a)=(1+2a)/4\) and
\(\Pr(Y=1\mid A=a,M=m)=(1+a+2m)/8\). The total-effect contrast is
\(1/4\). Holding the marginal distribution of \(M\) fixed while
standardizing the conditional contrasts gives \(1/8\). Thus a descendant
can change a total-effect calculation even when positivity holds. This
second contrast is a different functional of the stipulated law; it is
not a substitute for the total effect simply because both are numbers.
The wrapper's documented total/direct-effect switch must be an explicit
part of the request. A new claim about a direct-effect estimand would
need its own assumptions and interpretation. \citep{exec18_gp_graph_reference,exec18_gp_dagitty}

\subsection{Choose a runtime boundary without erasing graph uncertainty}

The wrapper documents graph classes beyond observed DAGs, but the proposed
first admission does not inherit every advertised class. The inspected
JavaScript validation code leaves proper MAG/PAG validation as a TODO, and
the inspected \texttt{pagToPdag} transformation changes endpoint marks and
graph type. Replacing a partial ancestral graph by a directed representation
may discard the ambiguity relevant to an intervention. A future extension
must preserve endpoint marks, latent and selection meanings, class validity
and the scope of the adjustment criterion. Syntax acceptance alone does not
supply those guarantees. Restricting the initial adapter to an observed DAG
keeps this unresolved composition visible. \citep{exec18_gp_dagitty}

Original-language reuse is preferred to an immediate port because it gives
a concrete comparison surface: parse the same supplied graph, validate the
same query, and compare typed adjustment results against independent finite
references. This is an engineering rationale, not evidence that a wrapper
already conforms. A direct JavaScript integration remains an alternative
when its preconditions and packaging are independently qualified; a port
would additionally need algorithmic equivalence, graph-class preservation
and a maintenance plan. Browser editing can support human review, but it
does not establish that a discovered or expert-supplied diagram is causal.
No browser, R, Node or upstream test execution supports the current choice.
\citep{exec18_gp_dagitty}

Rights also follow the selected components. The root contains GPL version 2
text; the R description declares GPL-2; one inspected helper has a
GPL version 2-or-later header, and other helpers declare MIT terms. The tree
contains vendored Underscore versions 1.8.3 and 1.13.1 and a separate test
dependency range. Those observations do not establish which copy a proposed
runtime loads, a vulnerability, or a complete transitive-rights inventory.
The release artifact, dependency resolution, notices, intended distribution
and source obligations must be examined together. Neither a separate process
nor a public download alone settles the rights of a particular arrangement.
The defensible proposal is pinned reference reuse with executable admission
deferred until that bounded trial and rights review are complete.
\citep{exec18_gp_dagitty}
